\section{Introduction}
\label{sec:intro}

Many training objectives are designed around a mechanism: they are meant to
improve a model by changing a named property of what it produces.
Frequency-aware losses for diffusion models are a clear case. Generated images
deviate from real ones in their high-frequency spectra
\citep{durall2020watch,dzanic2020fourier,frank2020leveraging,schwarz2021frequency},
and diffusion models in particular under-represent high spatial frequencies
\citep{ricker2024towards,yang2023diffusion}. Objectives that reweight the
diffusion loss across Fourier or wavelet components aim to correct this
deficit: by the current spectral error \citep{jiang2021focal}, inversely to
the spectrum magnitude \citep{yang2023diffusion}, on a schedule tied to the
noise level \citep{thao2026awwl}, or through auxiliary spectral terms
\citep{chandran2026spectral}.

Such objectives are evaluated by training them and comparing a metric with
that of a baseline, which shows whether an objective helps but not whether
its mechanism accounts for the result. An objective can change the statistic
it names while its effect on the metric arises through other channels, and
ablations vary the strength of a mechanism without measuring its value in the
metric; failing to identify the source of an empirical gain is a recurring
weakness of machine-learning research \citep{lipton2019troubling}. For
objectives motivated by a statistic of their output, the source can be
measured directly.

\paragraph{The pricing test.}
We propose a test that does so (Section~\ref{sec:pricing-test}). Let $B$
and $A$ be samples of a baseline model and of a model trained with the
objective under test, $M$ a metric (lower is better) and $s$ the statistic
that the objective claims to change. The test measures the change
$\delta_A = s(A) - s(B)$ that training achieved, imposes it on the baseline's
own samples with an operator $T_{\delta_A}$, and scores the result
(Fig.~\ref{fig:teaser}(a)). The value $V = M(B) - M(T_{\delta_A}(B))$ is
what the change is worth in $M$, and the mechanism residual gap
$\mathrm{MRG} = M(A) - M(T_{\delta_A}(B))$, a signed quantity, is the
metric of the trained objective minus that of the baseline carrying the same
change; with $D = M(A) - M(B)$ the observed gap, $\mathrm{MRG} = D + V$
exactly. A positive MRG means that what the objective alters besides $s$
costs the metric on balance, which does not imply $D > 0$; a negative MRG
means that the objective gains more than its change in $s$ is worth. The
split is analogous to the separation of direct and indirect effects in
mediation analysis \citep{robins1992identifiability,pearl2001direct}
(Section~\ref{sec:pricing-reading}). Its central identifying assumption is
that an operator imposing a change in $s$ moves $M$ as much as training that
produced the same change would, and six preconditions
(Section~\ref{sec:pricing-pre}) make its terms interpretable, among them that
the metric registers $s$ and that filtered images are scored without
re-quantisation. The value of a given change needs only the baseline's
samples; the MRG adds the objective, trained at the seeds of the comparison.

\guardedfigure{1.0}{teaser.png}{%
  The pricing test and its verdict on CIFAR-10. (a)~The change $\delta_A$ in
  the spectral statistic $s$ between the objective's samples $A$ and the
  baseline's samples $B$ is imposed on $B$ by the operator
  $T_A = T_{\delta_A}$, giving the value $V$ and the MRG, a signed difference.
  (b)~FID relative to $M(B)$, paired by training seed (lower is better), with
  paired 95\% $t$-intervals. Left: AWWL against MSE (five seeds, epoch 199,
  DDIM-100); the MRG has raw $p = \mrgMainP$ and Holm $p = \mrgArmAWWLHolm$
  over the \nFamilyPricing{} priced configurations. Right: the positive
  control, MSE with weight averaging at epoch 399 of an 800-epoch run, its
  FID optimum ($A$), against its epoch-199 checkpoint ($B$; three seeds;
  Section~\ref{sec:worth-checks}).}{fig:teaser}

Two controls check the test (Section~\ref{sec:worth-checks}): longer
training, which targets nothing spectral, returns a significantly negative
MRG, and a static band weighting that barely changes the spectrum returns a
value near zero and an MRG close to its FID gap. As $|V|$ is small in both,
$\mathrm{MRG} \approx D$, so the controls check the direction of the verdict
and the accounting, not the identifying assumption. A third control tests
that assumption directly: MSE trained on images that carry AWWL's spectral
change acquires the change through training and returns an MRG of
$\mrgPlanted$ $[\mrgPlantedLo, +\mrgPlantedHi]$, as the test requires when the
assumption holds (Section~\ref{sec:worth-planted}).

\paragraph{What the test finds.}
We apply the test to Adaptive Weight Wavelet Loss (AWWL)
\citep{thao2026awwl}, a wavelet band weighting that depends on the noise
level, introduced in our conference paper. The metric is the Fr\'echet Inception
Distance (FID) \citep{heusel2017gans}, the baseline the mean-squared-error
(MSE) objective, and each trained configuration (an arm) has five seeds
(three for four arms) on CIFAR-10, sampled with 100-step DDIM
\citep{song2021denoising}. No arm has a significantly lower FID than MSE
(Section~\ref{sec:audit}). AWWL's spectral change is worth
$V = \priceTransplantValue$~FID, and the value sits at mid radial
frequencies, the lower part of the range that AWWL up-weights, not at the
highest, where the baseline's deficit is largest. AWWL trails the baseline
carrying that change by $\mathrm{MRG} = \mrgMain$~FID (raw $p = \mrgMainP$),
a result that does not survive Holm correction across the \nFamilyPricing{}
priced arms ($p = \mrgArmAWWLHolm$; the verdict depends on the family), and its net gap to MSE is not resolved
(Section~\ref{sec:worth}).
Trained without the other, AWWL's timestep schedule and frequency shape at
constant total weight are each significantly worse than MSE, by
\dTimestep{} and \dSumFixed~FID, and partly cancel when combined, so how
AWWL's cost divides between them depends on the order; the shape produces
most of the spectral change, and fixing the total weight leaves the
effective timestep weight free (Section~\ref{sec:confound}). Across arms,
corrected MRG and FID-gap verdicts coincide, and every priced arm worse than
MSE, frequency-aware or not, has a positive MRG; the test adds the
decomposition, what each change is worth, not a more powerful test of the
gap (Section~\ref{sec:worth-arms}). Repeated in DINOv2 features, the test
resolves AWWL's gap to MSE and finds its spectral change worth nothing, so the
verdict does not rest on Inception features (Section~\ref{sec:worth-dino}).

Two results bear on evaluation beyond this objective. White noise within
$\pm0.5$ least significant bit (LSB) lowers the baseline's FID by
\priceNoiseHalfValue, so re-quantising filtered samples to 8~bits inflates
their value; on held-out real images the same noise does not lower FID, at
$\pm1$~LSB it raises it, and in DINOv2 features it raises the distance
(Sections~\ref{sec:worth-curve} and~\ref{sec:worth-dino}). On CelebA-64,
whose baseline has as large a high-band deficit as CIFAR-10's, spectral
corrections applied as a fixed radial gain raise FID; the two settings differ
in data, resolution, operating point and training steps, which the
experiment does not separate (Section~\ref{sec:celeba}).

\paragraph{Contributions.}
\begin{enumerate}
  \item A pricing test: what a change in an output statistic is worth, and
    where, from baseline samples alone, and an exact split of an objective's
    gap into that value and a residual (Section~\ref{sec:pricing-test}).
  \item A five-seed audit of AWWL: its conference results do not replicate,
    its printed equation differs from its code, and no objective improves
    significantly on MSE in FID (Section~\ref{sec:audit}).
  \item The test across \nFamilyPricing{} arms, separating what each spectral
    change is worth, and a two-by-two design crossing AWWL's timestep
    schedule with its frequency shape (Sections~\ref{sec:worth}
    and~\ref{sec:confound}).
  \item The limits of the test: a price that changes sign on CelebA-64,
    metrics that do not register the statistic in DreamBooth, and
    re-quantisation as a hazard (Sections~\ref{sec:worth-curve}
    and~\ref{sec:generality}).
\end{enumerate}

\paragraph{Relation to the conference version.}
AWWL was introduced by the present authors, with two further co-authors, in the proceedings of CITA 2026
\citep{thao2026awwl}, which reported single-run gains on CIFAR-10 and in
DreamBooth fine-tuning. The present paper is not an extended version of that
paper: its subject is the pricing test, which the conference paper does not
contain, and AWWL is the case to which the test is applied. No experiment,
table or figure of the conference paper is reused; the overlap is confined to
the definition of AWWL (Section~\ref{sec:family-awwl}) and to the published
results, which Table~\ref{tab:replication} quotes in order to test them.

Section~\ref{sec:related} reviews related work and
Section~\ref{sec:background} the objective family;
Sections~\ref{sec:pricing-test} and~\ref{sec:protocol} define the test and
the setup, Sections~\ref{sec:audit} to~\ref{sec:confound} audit AWWL and
price its change, Section~\ref{sec:generality} examines the reach of the test
and Section~\ref{sec:discussion} discusses the results and their limits.
